\section{From independent rows to action laws}\label{sec:actionlaws}

The pure-row theorem treats each named operator independently except for the shared carrier quotient.  Genuine action laws connect the row copies.  The raw-row formula and its feedback operator do not automatically
compute the forced quotient of these stronger presentations. In this
section $\theta_E^A$ always means equality in the full indicated ground
theory. We prove the structural consequences directly from that theory.

\subsection{Empty data produce a free action forest}

Assume $\Gamma=\varnothing$ and no action law beyond carrier compatibility.  For a word $w=b_1\cdots b_n\in B^*$ define
\[
W_w(a)
=
\alpha(b_1,\alpha(b_2,\ldots,\alpha(b_n,a)\ldots)),
\qquad
W_\varepsilon(a)=a.
\]

\begin{theorem}[Free action forest]\label{thm:forest}
In the empty-row typed theory, the map
\[
A\times B^*\longrightarrow \mathbb I_E^{\A},
\qquad
(a,w)\longmapsto[W_w(a)]
\]
is injective.  Consequently the initial carrier sort is infinite, and the number of distinct clean chains of length at most $d$ is at least
\[
|A|\sum_{k=0}^d|B|^k.
\]
\end{theorem}

\begin{proof}
Construct a model with operator sort $B$ and carrier sort $A\times B^*$.  Interpret carrier names by $(a,\varepsilon)$.  If a native $k$-ary carrier operation is applied to elements with a common word label $w$, put
\[
f((a_1,w),\ldots,(a_k,w))=(f^A(a_1,\ldots,a_k),w),
\]
and assign arbitrary values on mixed labels.  Define
\[
\alpha(b,(a,w))=(a,bw).
\]
The native diagrams hold at the base fiber and carrier compatibility holds because both sides of \eqref{eq:compat} have the same word label $b$.  Evaluation gives
\[
W_w(a)\mapsto(a,w),
\]
so distinct pairs are separated in a model of $E$ and hence cannot be equal in the initial algebra.
\end{proof}

The theorem identifies an embedded family of clean action chains.
It does not assert that these chains exhaust the initial carrier sort:
native operations applied to different fibers can produce further terms.

\begin{proposition}[Finiteness boundary for independent rows]\label{prop:shallowfinite}
In the pure typed row fragment, the initial carrier sort is finite if and only if every named ground action term $\alpha(b,a)$ is carrier-valued.
\end{proposition}

\begin{proof}
If all named action terms reduce to carrier names, structural induction reduces every closed carrier-sort term to a named carrier constant: native operations use $D_A$, and an action term first reduces its carrier argument to a name and then uses ground-action closure.

Conversely, suppose $x=\alpha(b,a)$ is external. The separation model
in Theorem~\ref{thm:globalrows} realizes $x$ as an element outside the
central named carrier. Choose a named $b_0$ and adjoin infinitely many
fresh elements $x_1,x_2,\ldots$. Preserve all native operations on the old
carrier and extend them arbitrarily to tuples using new elements.
Redefine $\alpha(b_0,x)=x_1$ and put
$\alpha(b_0,x_i)=x_{i+1}$; preserve the action on all central inputs.
Every pure-row axiom involves the action only on named central inputs
(possibly a native term equal to such an input), so these changes preserve
all axioms. The terms $x,\alpha(b_0,x),\alpha(b_0,\alpha(b_0,x)),\ldots$
have pairwise distinct values in the new model. Hence the initial carrier
sort is infinite.
\end{proof}

\subsection{Composition folds the forest}

Suppose now that $B$ is a finite semigroup with its named multiplication table and add the ground composition scheme
\begin{equation}\label{eq:Comp}
\mathrm{Comp}:\qquad
\alpha(gh,a)=\alpha(g,\alpha(h,a))
\end{equation}
for named $g,h\in B$ and named $a\in A$.

\begin{proposition}[Cayley folding]\label{prop:cayley}
In the empty-row theory with ground Comp, every nonempty clean chain satisfies
\[
E\vdash W_{b_1\cdots b_n}(a)
=
\alpha(\overline{b_1\cdots b_n},a),
\]
where the bar denotes the product in $B$.  Moreover, if no other action law is present, then
\[
E\vdash W_u(a)=W_v(a')
\quad\Longleftrightarrow\quad
\bigl(a=a'\text{ in }A\bigr)
\ \text{and}\
\bigl(\overline u=\overline v\text{ in }B\bigr)
\]
for nonempty words $u,v$.
\end{proposition}

\begin{proof}
Repeated bottom-up use of \eqref{eq:Comp} and the multiplication table reduces every clean chain to a single row term.  For the converse, use the carrier-sort model consisting of a base copy $A_*$ and copies $A_h$ indexed by $h\in B$, with
\[
\alpha(g,a_*)=a_g,
\qquad
\alpha(g,a_h)=a_{gh},
\]
and carrier operations acting copywise, with arbitrary values on mixed-copy tuples. Ground Comp and compatibility hold by construction, and different base elements and different product labels remain distinct.
\end{proof}

Thus Comp replaces the free word-indexed forest by a finite Cayley-type diagram of carrier copies
\[
A_h\xrightarrow{g}A_{gh}.
\]
This is the natural bridge from independent incomplete rows to partial-action and globalization theory.

\subsection{Units, inverses, and partial group actions}

When $B$ is a group, add the ground schemes
\[
\mathrm{Un}:\quad\alpha(e,a)=a,
\]
\[
\mathrm{Inv}:\quad\alpha(g^{-1},\alpha(g,a))=a,
\]
in addition to Comp.


For the theory with these additional laws, let $Q=A/\theta_E^A$ and define
$D_g$ semantically as the set of $[a]\in Q$ for which $\alpha(g,a)$ equals
a carrier name in the full initial algebra. The resulting forced value
map is $\tau_g:D_g\to Q$. These definitions are well defined by congruence;
we do not compute them by the pure-row saturation formula.

\begin{proposition}[Canonical partial group action]\label{prop:partialgroup}
Under ground Un, Comp, and Inv:
\begin{enumerate}[label=(\roman*)]
\item $D_g$ is closed under all positive-arity native operations. If it is
nonempty, it is a subalgebra, and $\tau_g$ is a homomorphism on its domain.
\item $D_e=Q$ and $\tau_e=\id_Q$.
\item $\tau_g:D_g\to D_{g^{-1}}$ is a bijection preserving native
operations in both directions, with inverse $\tau_{g^{-1}}$.
\item The exact domain of the composite is
\[
\operatorname{dom}(\tau_g\circ\tau_h)=D_h\cap D_{gh},
\]
and $\tau_g(\tau_h(x))=\tau_{gh}(x)$ on this domain.
\end{enumerate}
Thus the forced maps form a partial group action, allowing empty domains.
If every $D_g=Q$, they define a homomorphism $B\to\Aut(Q)$.
\end{proposition}

\begin{proof}
If each $x_i=[a_i]$ belongs to $D_g$, choose names $c_i$ for their forced
images. Ground compatibility gives
$\alpha(g,f(\bar a))=f(\bar c)$, which is a named carrier value by $D_A$.
This proves domain closure and preservation of operations. Un gives (ii).
If $\tau_g([a])=[c]$, substitution in ground Inv at $a$ gives
$\alpha(g^{-1},c)=a$. Applying the same argument with $g^{-1}$ proves both
inverse identities, establishing (iii).

If $x\in D_h$ and $\tau_h(x)\in D_g$, ground Comp identifies their
composite carrier value with $\alpha(gh,x)$, so $x\in D_{gh}$.
Conversely, if $x\in D_h\cap D_{gh}$, choose carrier names for
$\tau_h(x)$ and $\tau_{gh}(x)$. The same Comp instance states that applying
$g$ to the former equals the latter. Hence $\tau_h(x)\in D_g$.
This proves the domain identity and the value identity in (iv).
Full domains now give automorphisms satisfying the group law.
\end{proof}

Empty domains really occur: take a one-element carrier, $B=C_2$, and no
supplied cells. A two-point regular $C_2$-orbit, with the named carrier at
one point and native operations acting on constant fibers, is a model
whose nonidentity action value is external. Hence the corresponding
forced domain is empty. Positive arities make its closure condition
vacuous; no empty algebra is being required in our category.

This construction connects the forced maps to partial actions and
premorphisms~\cite{Hollings2007,KhrypchenkoNovikov2018}. Here the partial
maps are consequences of arbitrary ground prescriptions. Their carrier
quotient must be taken in the full theory including the group laws.
